% !TeX program = lualatex
% ================================================================
% Urdu Parallel Text Version of FactorisingAlgebraicFractionsQuickQuestions
%
% Generated by the server using the following settings:
%
%   language            Urdu (urd)
%   text direction      right to left
%   font                Noto Nastaliq Urdu
%   fallback font       Noto Serif
%   built from          latex/quickQuestions/factorisingAlgebraicFractionsQuickQuestions.tex
%   vocabulary key      latex/quickQuestions/factorisingAlgebraicFractionsQuickQuestions/L2/urd/factorisingAlgebraicFractionsQuickQuestions_urduKey.tex
%   tier 3 vocabulary   green   RGB 0 128 0
%   English text        black   RGB 0 0 0
%   Urdu text           purple  RGB 112 48 160
%   layout macros       version 3
%
% Please don't edit any information in this comment block - the server
% compares it against the current settings and rebuilds this file when
% they differ, so changes here would be overwritten. However, feel free
% to make updates to the LaTeX below if you want to edit it.
% ================================================================

\documentclass[aspectratio=169]{beamer}
% ================================================================
% Copied from the English sheet this was translated from, so that
% anything it sets up for itself still works here
% ================================================================
\usepackage{amsmath}
\usepackage{amssymb}
\usepackage{cancel}
\usepackage{tikz}
\usetikzlibrary{positioning}

% ================================================================
% Quick Questions deck: factorising and algebraic fractions.
%
% Each topic opens with five true/false statements and then runs ten
% quick questions that get gradually harder. Every question gets two
% slides: the question on its own for the mini whiteboards, then the
% same question again with the solution underneath in red.
%
% The two green topics are factorising itself, first into a single
% bracket and then into a double bracket. The two amber topics use
% that factorising to simplify an algebraic fraction, again single
% brackets before double ones. The two red topics do arithmetic with
% algebraic fractions: adding and subtracting, then multiplying and
% dividing.
%
% Where the working can be drawn it is: area rectangles for taking a
% common factor out, the 2 by 2 grid for the double brackets, a
% struck-through common factor for the cancelling, and a flow chart
% for dividing, where turning the second fraction upside down is a
% step of its own. Adding and subtracting is left as written working:
% an algebraic numerator cannot be shaded honestly, so a picture there
% would show a different quantity from the answer.
% ================================================================

\setbeamertemplate{navigation symbols}{}

% A link in the bottom left of every slide back to the contents, so a
% deck this long can be jumped around in during a lesson.
\setbeamertemplate{footline}{%
  \hbox to\paperwidth{%
    \hspace{1em}%
    \hyperlink{contents}{%
      \color{gray}\scriptsize$\leftarrow$~Back to contents}%
    \hfil}%
  \vskip4pt%
}

% Topic divider.
\newcommand{\qtopic}[1]{%
  \section{#1}%
  \begin{frame}
    \vfill
    \begin{center}
      {\usebeamercolor[fg]{structure}\Huge\bfseries #1}
    \end{center}
    \vfill
  \end{frame}%
}

% \qq[<question size>][<solution size>]{<question>}{<solution>}
% The two optional sizes drop the type for the wordier questions and
% the longer solutions, so nothing runs off the slide.
\NewDocumentCommand{\qq}{O{\Huge}O{\LARGE}+m+m}{%
  \begin{frame}
    \vfill
    \begin{center}
      \begin{minipage}{0.92\textwidth}
        \centering #1 #3
      \end{minipage}
    \end{center}
    \vfill
  \end{frame}%
  \begin{frame}
    \vfill
    \begin{center}
      \begin{minipage}{0.92\textwidth}
        \centering #1 #3
        \par\vspace{0.9em}
        {\color{red}#2 #4}
      \end{minipage}
    \end{center}
    \vfill
  \end{frame}%
}

% \tf[<statement size>][<answer size>]{<statement>}{<answer>}
% As \qq, but with a standing "True or false?" heading above the
% statement. The statements hunt for the usual misconceptions, so the
% answer always says why, not just which.
\NewDocumentCommand{\tf}{O{\LARGE}O{\Large}+m+m}{%
  \begin{frame}
    \vfill
    \begin{center}
      \begin{minipage}{0.92\textwidth}
        \centering
        {\usebeamercolor[fg]{structure}\Large\bfseries True or false?}
        \par\vspace{1em}
        #1 #3
      \end{minipage}
    \end{center}
    \vfill
  \end{frame}%
  \begin{frame}
    \vfill
    \begin{center}
      \begin{minipage}{0.92\textwidth}
        \centering
        {\usebeamercolor[fg]{structure}\Large\bfseries True or false?}
        \par\vspace{1em}
        #1 #3
        \par\vspace{0.9em}
        {\color{red}#2 #4}
      \end{minipage}
    \end{center}
    \vfill
  \end{frame}%
}

% Chained working, one equation per row, so that a long simplification
% does not have to be squeezed onto a single line. \dfrac is taller
% than the row strut, so rows carry an explicit skip such as \\[0.95em].
\newcommand{\work}[1]{%
  \begin{tabular}{@{}r@{\;=\;}l@{}}#1\end{tabular}%
}

% \strike{<expression>} --- a factor crossed out because it has been
% divided from the top and the bottom of a fraction.
\newcommand{\strike}[1]{\cancel{#1}}

% ================================================================
% Area pictures for factorising
% ================================================================

% \pairarea{<height>}{<width one>}{<width two>}{<area one>}{<area two>}
% One rectangle of height <height> cut into two parts. The two parts
% have the areas of the two terms being factorised, and the height is
% the factor taken out, so the whole rectangle is the factorised form.
\newcommand{\pairarea}[5]{%
  \begin{tikzpicture}[x=1cm,y=1cm,font=\small]
    \draw[red!60!black,thick,fill=red!8] (0,0) rectangle (3.2,1.5);
    \draw[red!60!black,thick,fill=red!8] (3.2,0) rectangle (5.0,1.5);
    \node[red] at (1.6,0.75) {$#4$};
    \node[red] at (4.1,0.75) {$#5$};
    \node[red,above=2pt] at (1.6,1.5) {$#2$};
    \node[red,above=2pt] at (4.1,1.5) {$#3$};
    \node[red,left=3pt] at (0,0.75) {$#1$};
  \end{tikzpicture}%
}

% \triplearea{<height>}{<w1>}{<w2>}{<w3>}{<a1>}{<a2>}{<a3>}
% The same picture cut into three parts, for a three term expression.
\newcommand{\triplearea}[7]{%
  \begin{tikzpicture}[x=1cm,y=1cm,font=\small]
    \draw[red!60!black,thick,fill=red!8] (0,0) rectangle (2.6,1.5);
    \draw[red!60!black,thick,fill=red!8] (2.6,0) rectangle (4.8,1.5);
    \draw[red!60!black,thick,fill=red!8] (4.8,0) rectangle (6.4,1.5);
    \node[red] at (1.3,0.75) {$#5$};
    \node[red] at (3.7,0.75) {$#6$};
    \node[red] at (5.6,0.75) {$#7$};
    \node[red,above=2pt] at (1.3,1.5) {$#2$};
    \node[red,above=2pt] at (3.7,1.5) {$#3$};
    \node[red,above=2pt] at (5.6,1.5) {$#4$};
    \node[red,left=3pt] at (0,0.75) {$#1$};
  \end{tikzpicture}%
}

% \quadgrid{<w1>}{<w2>}{<h1>}{<h2>}{<a>}{<b>}{<c>}{<d>}
% The 2 by 2 grid for a double bracket. The two widths are the terms
% of one bracket and the two heights the terms of the other, so the
% four cells add to the quadratic and the labels round the outside
% read off the factorised form.
\newcommand{\quadgrid}[8]{%
  \begin{tikzpicture}[x=1cm,y=1cm,font=\small]
    \draw[red!60!black,thick,fill=red!8] (0,0) rectangle (2.8,1.4);
    \draw[red!60!black,thick,fill=red!8] (2.8,0) rectangle (4.6,1.4);
    \draw[red!60!black,thick,fill=red!8] (0,-1.1) rectangle (2.8,0);
    \draw[red!60!black,thick,fill=red!8] (2.8,-1.1) rectangle (4.6,0);
    \node[red] at (1.4,0.7) {$#5$};
    \node[red] at (3.7,0.7) {$#6$};
    \node[red] at (1.4,-0.55) {$#7$};
    \node[red] at (3.7,-0.55) {$#8$};
    \node[red,above=2pt] at (1.4,1.4) {$#1$};
    \node[red,above=2pt] at (3.7,1.4) {$#2$};
    \node[red,left=3pt] at (0,0.7) {$#3$};
    \node[red,left=3pt] at (0,-0.55) {$#4$};
  \end{tikzpicture}%
}

% \flow{<start>}{<first step>}{<middle>}{<second step>}{<end>}
% A flow chart of three boxes joined by two labelled arrows. Used for
% dividing, where turning the second fraction upside down and then
% multiplying are separate steps, done one after the other. The step
% labels wrap to a fixed width, so a label of a few words sits above
% its own arrow instead of running into the boxes either side.
\newcommand{\flow}[5]{%
  \begin{tikzpicture}[
      fbox/.style={draw=red!60!black, rounded corners=2pt, fill=red!8,
                   inner sep=5pt, text=red, font=\Large},
      farr/.style={->, red, thick, shorten >=2pt, shorten <=2pt},
      flab/.style={above, midway, text=red, font=\small,
                   align=center, text width=3.2cm, inner sep=2pt,
                   execute at begin node={\hyphenpenalty=10000
                                          \exhyphenpenalty=10000}}]
    \node[fbox] (a) {$#1$};
    \node[fbox, right=3.4cm of a] (b) {$#3$};
    \node[fbox, right=3.4cm of b] (c) {$#5$};
    \draw[farr] (a) -- node[flab]{#2} (b);
    \draw[farr] (b) -- node[flab]{#4} (c);
  \end{tikzpicture}%
}

% Contents-slide bullets, colour-coded by difficulty band: green for
% the two factorising topics, orange for the two simplifying topics
% and red for the two arithmetic ones. The green is darkened because a
% pure green washes out on a projector.
\setbeamertemplate{section in toc}{%
  \ifnum\inserttocsectionnumber<3
    \textcolor{green!55!black}{$\bullet$}%
  \else
    \ifnum\inserttocsectionnumber<5
      \textcolor{orange}{$\bullet$}%
    \else
      \textcolor{red}{$\bullet$}%
    \fi
  \fi
  ~\inserttocsection%
}

\title{Factorising and Algebraic Fractions: Quick Questions}
\subtitle{Mini whiteboard questions, with solutions}
\author{}
\date{}

% Characters this language's font has not got are borrowed from another,
% rather than being silently dropped - see docs/translated-sheet-latex.md
\directlua{%
  if luaotfload and luaotfload.add_fallback then
    luaotfload.add_fallback("eallatinfallback",
      { "Noto Serif:mode=harf;" })
  else
    tex.error("this luaotfload is too old to register a fallback font")
  end
}
\usepackage[english,bidi=basic]{babel}
\babelprovide[import]{urdu}
\babelfont[urdu]{rm}[Renderer=Harfbuzz, BoldFont={Noto Nastaliq Urdu}, ItalicFont={Noto Nastaliq Urdu}, BoldItalicFont={Noto Nastaliq Urdu}, RawFeature={fallback=eallatinfallback}]{Noto Nastaliq Urdu}
\babelfont[urdu]{sf}[Renderer=Harfbuzz, BoldFont={Noto Nastaliq Urdu}, ItalicFont={Noto Nastaliq Urdu}, BoldItalicFont={Noto Nastaliq Urdu}, RawFeature={fallback=eallatinfallback}]{Noto Nastaliq Urdu}
\newcommand{\ealtext}[1]{\foreignlanguage{urdu}{#1}}
\newcommand{\ealtextblock}[1]{{\raggedleft\foreignlanguage{urdu}{#1}\par}}
% ================================================================
% EAL layout helpers
% ================================================================
\usepackage{xcolor}

\definecolor{ealkeycolour}{RGB}{0,128,0}
\definecolor{ealenglish}{RGB}{0,0,0}
\definecolor{ealtwocolour}{RGB}{112,48,160}

\newsavebox{\ealboxen}
\newsavebox{\ealboxtr}
\newlength{\ealglosswd}
\newlength{\ealglossraise}
\setlength{\ealglossraise}{1.15em}

% a tier 3 word, in English and in translation alike
\newcommand{\ealkey}[1]{\textcolor{ealkeycolour}{#1}}
\newcommand{\ealkeytr}[1]{\textcolor{ealkeycolour}{\ealtext{#1}}}

% One translated word above one English word. Built from kernel
% primitives, never a tabular, which would break wherever the sheet
% loads array, booktabs or tabularx. The box is as wide as the wider of
% the two, so a long translation pushes its neighbours aside instead of
% overlapping them, and it claims its full height so a table row or a
% TikZ node makes room for it.
\newcommand{\ealgloss}[2]{%
  \sbox{\ealboxen}{\ealkey{#1}}%
  \sbox{\ealboxtr}{\small\textcolor{ealtwocolour}{\ealtext{#2}}}%
  \setlength{\ealglosswd}{\wd\ealboxen}%
  \ifdim\wd\ealboxtr>\ealglosswd\setlength{\ealglosswd}{\wd\ealboxtr}\fi
  \leavevmode
  \makebox[\ealglosswd][c]{%
    \makebox[0pt][c]{\usebox{\ealboxen}}%
    \makebox[0pt][c]{\raisebox{\ealglossraise}{\usebox{\ealboxtr}}}%
  }%
}

% opens up line spacing around a block of glosses, so the raised
% translations sit clear of the line above
\newenvironment{ealglossed}{\par\linespread{2.1}\selectfont\ignorespaces}{\par}

% a whole translated sentence above its English counterpart. block level,
% so it does not belong inside a tabular cell
\newcommand{\ealpara}[2]{%
  \par\addvspace{0.5\baselineskip}%
  {\color{ealtwocolour}\ealtextblock{#1}}%
  \penalty200
  {\color{ealenglish}#2\par}%
  \addvspace{0.5\baselineskip}%
}

\begin{document}

\frame{\titlepage}

\begin{frame}[label=contents]{Contents}
  \small
  \tableofcontents
\end{frame}

%================================================================
\qtopic{\ealkey{Factorising} into single \ealkey{brackets}}
%================================================================

\tf[\LARGE][\large]{$6x+9=3(2x+3)$}
  {\ealpara{\textbf{صحیح۔} $6x$ اور $9$ کا \ealkeytr{سب سے بڑا مشترک عامل} $3$ ہے، اور $3\times 2x=6x$، $3\times 3=9$۔}{\textbf{True.} The \ealkey{highest common factor} of $6x$ and $9$ is $3$, and $3\times 2x=6x$, $3\times 3=9$.}
   \par\vspace{0.5em}
   \pairarea{3}{2x}{3}{6x}{9}}

\tf[\Large][\large]{\ealpara{$8x+12$ مکمل طور پر \ealkeytr{تجزیه} ہو کر $2(4x+6)$ ہے}{$8x+12$ is fully \ealkey{factorised} as $2(4x+6)$}}
  {\ealpara{\textbf{غلط۔} $2$، $8x$ اور $12$ کا \ealkeytr{مشترک عامل} ہے، لیکن یہ \ealkeytr{سب سے بڑا مشترک عامل} نہیں ہے، اس لیے \ealkeytr{اظہار} مکمل طور پر \ealkeytr{تجزیه} نہیں ہوا۔}{\textbf{False.} $2$ is a \ealkey{common factor} of $8x$ and $12$, but it is not the \ealkey{highest common factor}, so the \ealkey{expression} is not fully \ealkey{factorised}.}
   \ealpara{$8x$ اور $12$ کا \ealkeytr{سب سے بڑا مشترک عامل} $4$ ہے، اس لیے $8x+12=4(2x+3)$۔}{The \ealkey{highest common factor} of $8x$ and $12$ is $4$, so $8x+12=4(2x+3)$.}}

\tf{$5x+10y=5(x+10y)$}
  {\ealpara{\textbf{غلط۔} \ealkeytr{بریکٹ} کے اندر ہر \ealkeytr{رکن} کو \ealkeytr{مشترک عامل} سے \ealkeytr{تقسیم} کرنا چاہیے، صرف پہلا \ealkeytr{رکن} نہیں۔}{\textbf{False.} Every \ealkey{term} inside the \ealkey{bracket} must be \ealkey{divided} by the \ealkey{common factor}, not only the first \ealkey{term}.}
   \ealpara{$10y\div 5=2y$، اس لیے $5x+10y=5(x+2y)$۔}{$10y\div 5=2y$, so $5x+10y=5(x+2y)$.}}

\tf{$x^{2}+x=x(x)$}
  {\ealpara{\textbf{غلط۔} $x(x)$ تو $x^{2}$ ہے، اس لیے دوسرا \ealkeytr{رکن} گم ہو گیا ہے۔}{\textbf{False.} $x(x)$ is $x^{2}$, so the second \ealkey{term} has been lost.}
   \ealpara{$x\div x=1$، اس لیے $x^{2}+x=x(x+1)$۔}{$x\div x=1$, so $x^{2}+x=x(x+1)$.}}

\tf{$3x+3=3(x+1)$}
  {\ealpara{\textbf{صحیح۔} دوسرا \ealkeytr{رکن} مکمل طور پر \ealkeytr{مشترک عامل} ہے، اور $3\div 3=1$، اس لیے \ealkeytr{بریکٹ} کے اندر $1$ رہ جاتا ہے۔}{\textbf{True.} The second \ealkey{term} is entirely the \ealkey{common factor}, and $3\div 3=1$, so a $1$ is left inside the \ealkey{bracket}.}}

\qq[\Huge][\large]{\ealpara{\ealkeytr{تجزیه کریں} $2x+6$}{\ealkey{Factorise} $2x+6$}}
  {\ealpara{$2x$ اور $6$ کا \ealkeytr{سب سے بڑا مشترک عامل} $2$ ہے۔}{The \ealkey{highest common factor} of $2x$ and $6$ is $2$.}
   $2x+6=2(x+3)$
   \par\vspace{0.5em}
   \pairarea{2}{x}{3}{2x}{6}}

\qq[\Huge][\LARGE]{\ealpara{\ealkeytr{تجزیه کریں} $9x+24$}{\ealkey{Factorise} $9x+24$}}
  {\ealpara{$9x$ اور $24$ کا \ealkeytr{سب سے بڑا مشترک عامل} $3$ ہے۔}{The \ealkey{highest common factor} of $9x$ and $24$ is $3$.}
   \ealpara{$9x\div 3=3x$ اور $24\div 3=8$۔}{$9x\div 3=3x$ and $24\div 3=8$.}
   $9x+24=3(3x+8)$}

\qq[\Huge][\LARGE]{\ealpara{\ealkeytr{تجزیه کریں} $6y-15$}{\ealkey{Factorise} $6y-15$}}
  {\ealpara{$6y$ اور $15$ کا \ealkeytr{سب سے بڑا مشترک عامل} $3$ ہے۔}{The \ealkey{highest common factor} of $6y$ and $15$ is $3$.}
   $6y-15=3(2y-5)$}

\qq[\Huge][\LARGE]{\ealpara{\ealkeytr{تجزیه کریں} $10a+15b$}{\ealkey{Factorise} $10a+15b$}}
  {\ealpara{$10a$ اور $15b$ کا \ealkeytr{سب سے بڑا مشترک عامل} $5$ ہے۔}{The \ealkey{highest common factor} of $10a$ and $15b$ is $5$.}
   $10a+15b=5(2a+3b)$}

\qq[\Huge][\large]{\ealpara{\ealkeytr{تجزیه کریں} $x^{2}+7x$}{\ealkey{Factorise} $x^{2}+7x$}}
  {\ealpara{$x^{2}$ اور $7x$ کا \ealkeytr{سب سے بڑا مشترک عامل} $x$ ہے۔}{The \ealkey{highest common factor} of $x^{2}$ and $7x$ is $x$.}
   $x^{2}+7x=x(x+7)$
   \par\vspace{0.5em}
   \pairarea{x}{x}{7}{x^{2}}{7x}}

\qq[\Huge][\large]{\ealpara{\ealkeytr{تجزیه کریں} $12x^{2}+8x$}{\ealkey{Factorise} $12x^{2}+8x$}}
  {\ealpara{$12x^{2}$ اور $8x$ کا \ealkeytr{سب سے بڑا مشترک عامل} $4x$ ہے۔}{The \ealkey{highest common factor} of $12x^{2}$ and $8x$ is $4x$.}
   $12x^{2}+8x=4x(3x+2)$
   \par\vspace{0.5em}
   \pairarea{4x}{3x}{2}{12x^{2}}{8x}}

\qq[\Huge][\LARGE]{\ealpara{\ealkeytr{تجزیه کریں} $6x^{3}-9x^{2}$}{\ealkey{Factorise} $6x^{3}-9x^{2}$}}
  {\ealpara{اعداد کا \ealkeytr{سب سے بڑا مشترک عامل} $3$ ہے اور $x$ کی \ealkeytr{قوتوں} کا $x^{2}$ ہے، اس لیے \ealkeytr{سب سے بڑا مشترک عامل} $3x^{2}$ ہے۔}{The \ealkey{highest common factor} of the numbers is $3$ and of the \ealkey{powers} of $x$ is $x^{2}$, so the \ealkey{highest common factor} is $3x^{2}$.}
   $6x^{3}-9x^{2}=3x^{2}(2x-3)$}

\qq[\Huge][\LARGE]{\ealpara{\ealkeytr{تجزیه کریں} $4x^{2}y+6xy^{2}$}{\ealkey{Factorise} $4x^{2}y+6xy^{2}$}}
  {\ealpara{دونوں \ealkeytr{ارکان} میں $2$، $x$ اور $y$ موجود ہیں، اس لیے \ealkeytr{سب سے بڑا مشترک عامل} $2xy$ ہے۔}{Both \ealkey{terms} contain $2$, $x$ and $y$, so the \ealkey{highest common factor} is $2xy$.}
   $4x^{2}y+6xy^{2}=2xy(2x+3y)$}

\qq[\Huge][\large]{\ealpara{\ealkeytr{تجزیه کریں} $8x^{3}+12x^{2}+20x$}{\ealkey{Factorise} $8x^{3}+12x^{2}+20x$}}
  {\ealpara{$8x^{3}$، $12x^{2}$ اور $20x$ کا \ealkeytr{سب سے بڑا مشترک عامل} $4x$ ہے۔}{The \ealkey{highest common factor} of $8x^{3}$, $12x^{2}$ and $20x$ is $4x$.}
   $8x^{3}+12x^{2}+20x=4x(2x^{2}+3x+5)$
   \par\vspace{0.5em}
   \triplearea{4x}{2x^{2}}{3x}{5}{8x^{3}}{12x^{2}}{20x}}

\qq[\LARGE][\large]{\ealpara{\ealkeytr{تجزیه کریں} $5(x+1)+x(x+1)$}{\ealkey{Factorise} $5(x+1)+x(x+1)$}}
  {\ealpara{\ealkeytr{بریکٹ} $(x+1)$ دونوں \ealkeytr{ارکان} کا \ealkeytr{مشترک عامل} ہے۔}{The \ealkey{bracket} $(x+1)$ is a \ealkey{common factor} of both \ealkey{terms}.}
   $5(x+1)+x(x+1)=(x+1)(x+5)$
   \par\vspace{0.5em}
   \pairarea{x+1}{x}{5}{x(x+1)}{5(x+1)}}

%================================================================
\qtopic{\ealkey{Factorising} into double \ealkey{brackets}}
%================================================================

\tf[\LARGE][\large]{$x^{2}-9=(x-3)(x-3)$}
  {\ealpara{\textbf{غلط۔} $(x-3)(x-3)$ \ealkeytr{توسیع} ہو کر $x^{2}-6x+9$ بنتا ہے۔}{\textbf{False.} $(x-3)(x-3)$ \ealkey{expands} to $x^{2}-6x+9$.}
   \ealpara{\ealkeytr{علامات} مخالف ہونی چاہئیں: $x^{2}-9=(x+3)(x-3)$۔}{The \ealkey{signs} must be opposite: $x^{2}-9=(x+3)(x-3)$.}
   \par\vspace{0.4em}
   \quadgrid{x}{3}{x}{-3}{x^2}{3x}{-3x}{-9}}

\tf{$x^{2}+9=(x+3)(x+3)$}
  {\ealpara{\textbf{غلط۔} $(x+3)(x+3)$ \ealkeytr{توسیع} ہو کر $x^{2}+6x+9$ بنتا ہے۔}{\textbf{False.} $(x+3)(x+3)$ \ealkey{expands} to $x^{2}+6x+9$.}
   \ealpara{$x^{2}+9$ دو \ealkeytr{مربعوں} کا \ealkeytr{حاصل جمع} ہے، اور یہ \ealkeytr{تجزیه} نہیں ہوتا۔}{$x^{2}+9$ is a \ealkey{sum} of two \ealkey{squares}, and it does not \ealkey{factorise}.}}

\tf{$x^{2}+2x-15=(x+5)(x-3)$}
  {\ealpara{\textbf{صحیح۔} $5\times(-3)=-15$ اور $5+(-3)=2$۔}{\textbf{True.} $5\times(-3)=-15$ and $5+(-3)=2$.}}

\tf{$x^{2}-2x-8=(x+4)(x-2)$}
  {\ealpara{\textbf{غلط۔} $(x+4)(x-2)$ \ealkeytr{توسیع} ہو کر $x^{2}+2x-8$ بنتا ہے، اس لیے \ealkeytr{علامات} الٹی ہیں۔}{\textbf{False.} $(x+4)(x-2)$ \ealkey{expands} to $x^{2}+2x-8$, so the \ealkey{signs} are the wrong way round.}
   $x^{2}-2x-8=(x-4)(x+2)$}

\tf[\Large][\large]{\ealpara{$x^{2}+7x+12$ کو \ealkeytr{تجزیه} کرنے کے لیے ہمیں دو اعداد چاہیے جو $12$ کو \ealkeytr{ضرب} دیں اور $7$ میں \ealkeytr{جمع} ہوں}{To \ealkey{factorise} $x^{2}+7x+12$ we need two numbers that \ealkey{multiply} to $12$ and \ealkey{add} to $7$}}
  {\ealpara{\textbf{صحیح۔} دونوں اعداد \ealkeytr{مستقل حد} کو \ealkeytr{ضرب} دیتے ہیں اور $x$ کے \ealkeytr{عددی سر} میں \ealkeytr{جمع} ہوتے ہیں۔}{\textbf{True.} The two numbers \ealkey{multiply} to the \ealkey{constant term} and \ealkey{add} to the \ealkey{coefficient} of $x$.}
   \ealpara{اعداد $3$ اور $4$ ہیں، اس لیے $x^{2}+7x+12=(x+3)(x+4)$۔}{The numbers are $3$ and $4$, so $x^{2}+7x+12=(x+3)(x+4)$.}}

\qq[\Huge][\large]{\ealpara{\ealkeytr{تجزیه کریں} $x^{2}+5x+6$}{\ealkey{Factorise} $x^{2}+5x+6$}}
  {\ealpara{دو اعداد جو $6$ کو \ealkeytr{ضرب} دیں اور $5$ میں \ealkeytr{جمع} ہوں: $2$ اور $3$۔}{Two numbers that \ealkey{multiply} to $6$ and \ealkey{add} to $5$: $2$ and $3$.}
   $x^{2}+5x+6=(x+2)(x+3)$
   \par\vspace{0.5em}
   \quadgrid{x}{2}{x}{3}{x^2}{2x}{3x}{6}}

\qq[\Huge][\LARGE]{\ealpara{\ealkeytr{تجزیه کریں} $x^{2}+9x+20$}{\ealkey{Factorise} $x^{2}+9x+20$}}
  {\ealpara{دو اعداد جو $20$ کو \ealkeytr{ضرب} دیں اور $9$ میں \ealkeytr{جمع} ہوں: $4$ اور $5$۔}{Two numbers that \ealkey{multiply} to $20$ and \ealkey{add} to $9$: $4$ and $5$.}
   $x^{2}+9x+20=(x+4)(x+5)$}

\qq[\Huge][\large]{\ealpara{\ealkeytr{تجزیه کریں} $x^{2}+3x-10$}{\ealkey{Factorise} $x^{2}+3x-10$}}
  {\ealpara{دو اعداد جو $-10$ کو \ealkeytr{ضرب} دیں اور $3$ میں \ealkeytr{جمع} ہوں: $5$ اور $-2$۔}{Two numbers that \ealkey{multiply} to $-10$ and \ealkey{add} to $3$: $5$ and $-2$.}
   $x^{2}+3x-10=(x+5)(x-2)$
   \par\vspace{0.5em}
   \quadgrid{x}{5}{x}{-2}{x^2}{5x}{-2x}{-10}}

\qq[\Huge][\LARGE]{\ealpara{\ealkeytr{تجزیه کریں} $x^{2}-7x+12$}{\ealkey{Factorise} $x^{2}-7x+12$}}
  {\ealpara{\ealkeytr{مستقل حد} مثبت ہے اور $x$ کا \ealkeytr{عددی سر} منفی ہے، اس لیے دونوں اعداد منفی ہیں: $-3$ اور $-4$۔}{The \ealkey{constant term} is positive and the \ealkey{coefficient} of $x$ is negative, so both numbers are negative: $-3$ and $-4$.}
   $x^{2}-7x+12=(x-3)(x-4)$}

\qq[\Huge][\LARGE]{\ealpara{\ealkeytr{تجزیه کریں} $x^{2}-2x-15$}{\ealkey{Factorise} $x^{2}-2x-15$}}
  {\ealpara{دو اعداد جو $-15$ کو \ealkeytr{ضرب} دیں اور $-2$ میں \ealkeytr{جمع} ہوں: $-5$ اور $3$۔}{Two numbers that \ealkey{multiply} to $-15$ and \ealkey{add} to $-2$: $-5$ and $3$.}
   $x^{2}-2x-15=(x-5)(x+3)$}

\qq[\Huge][\LARGE]{\ealpara{\ealkeytr{تجزیه کریں} $x^{2}-49$}{\ealkey{Factorise} $x^{2}-49$}}
  {\ealpara{یہ \ealkeytr{دو مربعوں کا فرق} ہے، جہاں $x^{2}=(x)^{2}$ اور $49=7^{2}$۔}{This is a \ealkey{difference of two squares}, with $x^{2}=(x)^{2}$ and $49=7^{2}$.}
   $x^{2}-49=(x+7)(x-7)$}

\qq[\Huge][\large]{\ealpara{\ealkeytr{تجزیه کریں} $2x^{2}+7x+3$}{\ealkey{Factorise} $2x^{2}+7x+3$}}
  {\ealpara{$2x^{2}$، $2x\times x$ سے آتا ہے، اور $3$، $1\times 3$ سے۔}{The $2x^{2}$ comes from $2x\times x$, and the $3$ from $1\times 3$.}
   $2x^{2}+7x+3=(2x+1)(x+3)$
   \par\vspace{0.5em}
   \quadgrid{2x}{1}{x}{3}{2x^2}{x}{6x}{3}}

\qq[\Huge][\large]{\ealpara{\ealkeytr{تجزیه کریں} $2x^{2}+10x+12$}{\ealkey{Factorise} $2x^{2}+10x+12$}}
  {\ealpara{پہلے $2$ کا \ealkeytr{مشترک عامل} باہر نکالیں۔}{Take the \ealkey{common factor} of $2$ out first.}
   $2x^{2}+10x+12=2(x^{2}+5x+6)$
   \ealpara{$x^{2}+5x+6=(x+2)(x+3)$، اس لیے $2x^{2}+10x+12=2(x+2)(x+3)$۔}{$x^{2}+5x+6=(x+2)(x+3)$, so $2x^{2}+10x+12=2(x+2)(x+3)$.}
   \par\vspace{0.5em}
   \triplearea{2}{x^{2}}{5x}{6}{2x^{2}}{10x}{12}}

\qq[\Huge][\LARGE]{\ealpara{\ealkeytr{تجزیه کریں} $9x^{2}-16$}{\ealkey{Factorise} $9x^{2}-16$}}
  {\ealpara{یہ \ealkeytr{دو مربعوں کا فرق} ہے، جہاں $9x^{2}=(3x)^{2}$ اور $16=4^{2}$۔}{This is a \ealkey{difference of two squares}, with $9x^{2}=(3x)^{2}$ and $16=4^{2}$.}
   $9x^{2}-16=(3x+4)(3x-4)$}

\qq[\Huge][\large]{\ealpara{\ealkeytr{تجزیه کریں} $6x^{2}-x-12$}{\ealkey{Factorise} $6x^{2}-x-12$}}
  {\ealpara{$6x^{2}$، $3x\times 2x$ سے آتا ہے، اور $-12$، $4\times(-3)$ سے۔}{The $6x^{2}$ comes from $3x\times 2x$, and the $-12$ from $4\times(-3)$.}
   $6x^{2}-x-12=(3x+4)(2x-3)$
   \par\vspace{0.5em}
   \quadgrid{3x}{4}{2x}{-3}{6x^2}{8x}{-9x}{-12}}

%================================================================
\qtopic{\ealkey{Simplifying} \ealkey{algebraic} \ealkey{fractions} with single \ealkey{brackets}}
%================================================================

\tf[\LARGE][\large]{$\dfrac{x+2}{2}=x$}
  {\ealpara{\textbf{غلط۔} نیچے والا $2$ پورے \ealkeytr{صورت} کو \ealkeytr{تقسیم} کرتا ہے، صرف \ealkeytr{رکن} $2$ کو نہیں۔}{\textbf{False.} The $2$ on the bottom \ealkey{divides} the whole \ealkey{numerator}, not the \ealkey{term} $2$ on its own.}
   \ealpara{$x+2$ میں کوئی \ealkeytr{مشترک عامل} نہیں ہے، اس لیے $\dfrac{x+2}{2}$ \ealkeytr{سادہ} نہیں ہوتا۔}{$x+2$ has no \ealkey{common factor}, so $\dfrac{x+2}{2}$ does not \ealkey{simplify}.}
   \ealpara{جب $x=4$: $\dfrac{4+2}{2}=3$، نہ کہ $4$۔}{When $x=4$: $\dfrac{4+2}{2}=3$, not $4$.}}

\tf[\LARGE][\large]{$\dfrac{3x+6}{3}=x+2$}
  {\ealpara{\textbf{صحیح۔} $3$ پوری \ealkeytr{صورت} کا \ealkeytr{عامل} ہے: $3x+6=3(x+2)$۔}{\textbf{True.} $3$ is a \ealkey{factor} of the whole \ealkey{numerator}: $3x+6=3(x+2)$.}
   $\dfrac{\strike{3}(x+2)}{\strike{3}}=x+2$}

\tf[\LARGE][\large]{$\dfrac{x+3}{x}=3$}
  {\ealpara{\textbf{غلط۔} $x$ \ealkeytr{صورت} کا ایک \ealkeytr{رکن} ہے، پوری \ealkeytr{صورت} کا \ealkeytr{عامل} نہیں۔}{\textbf{False.} $x$ is a \ealkey{term} of the \ealkey{numerator}, not a \ealkey{factor} of the whole \ealkey{numerator}.}
   \ealpara{$\dfrac{x+3}{x}$ \ealkeytr{سادہ} نہیں ہوتا۔}{$\dfrac{x+3}{x}$ does not \ealkey{simplify}.}
   \ealpara{جب $x=1$: $\dfrac{1+3}{1}=4$، نہ کہ $3$۔}{When $x=1$: $\dfrac{1+3}{1}=4$, not $3$.}}

\tf[\LARGE][\large]{$\dfrac{6x+9}{3}=2x+9$}
  {\ealpara{\textbf{غلط۔} \ealkeytr{صورت} کے ہر \ealkeytr{رکن} کو $3$ سے \ealkeytr{تقسیم} کیا جاتا ہے، صرف پہلا \ealkeytr{رکن} نہیں۔}{\textbf{False.} Every \ealkey{term} of the \ealkey{numerator} is \ealkey{divided} by $3$, not the first \ealkey{term} on its own.}
   \ealpara{$6x+9=3(2x+3)$، اس لیے \ealkeytr{کسر} \ealkeytr{سادہ} ہو کر $2x+3$ بنتی ہے۔}{$6x+9=3(2x+3)$, so the \ealkey{fraction} \ealkey{simplifies} to $2x+3$.}}

\tf[\LARGE][\large]{$\dfrac{5x+10}{5x}=\dfrac{x+2}{x}$}
  {\ealpara{\textbf{صحیح۔} $5$ پوری \ealkeytr{صورت} اور پورے \ealkeytr{مخرج} کا \ealkeytr{مشترک عامل} ہے۔}{\textbf{True.} $5$ is a \ealkey{common factor} of the whole \ealkey{numerator} and the whole \ealkey{denominator}.}
   $\dfrac{\strike{5}(x+2)}{\strike{5}x}=\dfrac{x+2}{x}$}

\qq[\LARGE][\normalsize]{\ealpara{\ealkeytr{سادہ} کریں $\dfrac{2x+10}{2}$}{\ealkey{Simplify} $\dfrac{2x+10}{2}$}}
  {\ealpara{\ealkeytr{تجزیه} کریں \ealkeytr{صورت} کو۔}{\ealkey{Factorise} the \ealkey{numerator}.}
   $\dfrac{2x+10}{2}=\dfrac{\strike{2}(x+5)}{\strike{2}}=x+5$
   \par\vspace{0.5em}
   \pairarea{2}{x}{5}{2x}{10}}

\qq[\LARGE][\large]{\ealpara{\ealkeytr{سادہ} کریں $\dfrac{5x-15}{5}$}{\ealkey{Simplify} $\dfrac{5x-15}{5}$}}
  {\ealpara{\ealkeytr{صورت} کا \ealkeytr{سب سے بڑا مشترک عامل} $5$ ہے۔}{The \ealkey{highest common factor} of the \ealkey{numerator} is $5$.}
   $\dfrac{5(x-3)}{5}=\dfrac{\strike{5}(x-3)}{\strike{5}}=x-3$}

\qq[\LARGE][\normalsize]{\ealpara{\ealkeytr{سادہ} کریں $\dfrac{2x+8}{x+4}$}{\ealkey{Simplify} $\dfrac{2x+8}{x+4}$}}
  {\ealpara{\ealkeytr{تجزیه} کریں \ealkeytr{صورت} کو۔}{\ealkey{Factorise} the \ealkey{numerator}.}
   $\dfrac{2x+8}{x+4}=\dfrac{2\strike{(x+4)}}{\strike{(x+4)}}=2$
   \par\vspace{0.5em}
   \pairarea{2}{x}{4}{2x}{8}}

\qq[\LARGE][\large]{\ealpara{\ealkeytr{سادہ} کریں $\dfrac{3x-12}{x-4}$}{\ealkey{Simplify} $\dfrac{3x-12}{x-4}$}}
  {\ealpara{\ealkeytr{تجزیه} کریں \ealkeytr{صورت}: $3x-12=3(x-4)$۔}{\ealkey{Factorise} the \ealkey{numerator}: $3x-12=3(x-4)$.}
   $\dfrac{3\strike{(x-4)}}{\strike{(x-4)}}=3$}

\qq[\LARGE][\normalsize]{\ealpara{\ealkeytr{سادہ} کریں $\dfrac{x^{2}+3x}{x}$}{\ealkey{Simplify} $\dfrac{x^{2}+3x}{x}$}}
  {\ealpara{\ealkeytr{صورت} کا \ealkeytr{سب سے بڑا مشترک عامل} $x$ ہے۔}{The \ealkey{highest common factor} of the \ealkey{numerator} is $x$.}
   $\dfrac{x(x+3)}{x}=\dfrac{\strike{x}(x+3)}{\strike{x}}=x+3$
   \par\vspace{0.5em}
   \pairarea{x}{x}{3}{x^{2}}{3x}}

\qq[\LARGE][\normalsize]{\ealpara{\ealkeytr{سادہ} کریں $\dfrac{3x^{2}+6x+9}{3}$}{\ealkey{Simplify} $\dfrac{3x^{2}+6x+9}{3}$}}
  {\ealpara{$3$ \ealkeytr{صورت} کے تینوں \ealkeytr{ارکان} کا \ealkeytr{عامل} ہے۔}{$3$ is a \ealkey{factor} of all three \ealkey{terms} of the \ealkey{numerator}.}
   $\dfrac{\strike{3}(x^{2}+2x+3)}{\strike{3}}=x^{2}+2x+3$
   \par\vspace{0.5em}
   \triplearea{3}{x^{2}}{2x}{3}{3x^{2}}{6x}{9}}

\qq[\LARGE][\normalsize]{\ealpara{\ealkeytr{سادہ} کریں $\dfrac{4x^{2}+6x}{2x}$}{\ealkey{Simplify} $\dfrac{4x^{2}+6x}{2x}$}}
  {\ealpara{\ealkeytr{صورت} کا \ealkeytr{سب سے بڑا مشترک عامل} $2x$ ہے۔}{The \ealkey{highest common factor} of the \ealkey{numerator} is $2x$.}
   $\dfrac{2x(2x+3)}{2x}=\dfrac{\strike{2x}(2x+3)}{\strike{2x}}=2x+3$
   \par\vspace{0.5em}
   \pairarea{2x}{2x}{3}{4x^{2}}{6x}}

\qq[\LARGE][\large]{\ealpara{\ealkeytr{سادہ} کریں $\dfrac{4x+12}{6x+18}$}{\ealkey{Simplify} $\dfrac{4x+12}{6x+18}$}}
  {\ealpara{\ealkeytr{تجزیه} کریں \ealkeytr{صورت} اور \ealkeytr{مخرج} کو۔}{\ealkey{Factorise} the \ealkey{numerator} and the \ealkey{denominator}.}
   \par\vspace{0.4em}
   \work{$\dfrac{4x+12}{6x+18}$ & $\dfrac{4(x+3)}{6(x+3)}$\\[0.95em]
         $\dfrac{4\strike{(x+3)}}{6\strike{(x+3)}}$ & $\dfrac{4}{6}=\dfrac{2}{3}$}}

\qq[\LARGE][\large]{\ealpara{\ealkeytr{سادہ} کریں $\dfrac{10}{5x+20}$}{\ealkey{Simplify} $\dfrac{10}{5x+20}$}}
  {\ealpara{\ealkeytr{تجزیه} کریں \ealkeytr{مخرج}: $5x+20=5(x+4)$۔}{\ealkey{Factorise} the \ealkey{denominator}: $5x+20=5(x+4)$.}
   $\dfrac{\strike{5}\times 2}{\strike{5}(x+4)}=\dfrac{2}{x+4}$}

\qq[\LARGE][\large]{\ealpara{\ealkeytr{سادہ} کریں $\dfrac{6x^{2}+9x}{4x^{2}+6x}$}{\ealkey{Simplify} $\dfrac{6x^{2}+9x}{4x^{2}+6x}$}}
  {\ealpara{\ealkeytr{سب سے بڑا مشترک عامل} اوپر $3x$ اور نیچے $2x$ ہے۔}{The \ealkey{highest common factor} is $3x$ on the top and $2x$ on the bottom.}
   \par\vspace{0.4em}
   \work{$\dfrac{6x^{2}+9x}{4x^{2}+6x}$ & $\dfrac{3x(2x+3)}{2x(2x+3)}$\\[0.95em]
         $\dfrac{3\strike{x}\strike{(2x+3)}}{2\strike{x}\strike{(2x+3)}}$ & $\dfrac{3}{2}$}}

%================================================================
\qtopic{\ealkey{Simplifying} \ealkey{algebraic} \ealkey{fractions} with double \ealkey{brackets}}
%================================================================

\tf[\LARGE][\large]{$\dfrac{x^{2}+5x+6}{x+2}=x+3$}
  {\ealpara{\textbf{صحیح۔} $x^{2}+5x+6=(x+2)(x+3)$، اس لیے $(x+2)$ پوری \ealkeytr{صورت} کا \ealkeytr{عامل} ہے اور \ealkeytr{حذف} ہو جاتا ہے۔}{\textbf{True.} $x^{2}+5x+6=(x+2)(x+3)$, so $(x+2)$ is a \ealkey{factor} of the whole \ealkey{numerator} and \ealkey{cancels}.}
   $\dfrac{\strike{(x+2)}(x+3)}{\strike{(x+2)}}=x+3$}

\tf[\LARGE][\large]{$\dfrac{x^{2}+7x}{x^{2}+2x}=\dfrac{7}{2}$}
  {\ealpara{\textbf{غلط۔} $x^{2}$ ایک \ealkeytr{رکن} ہے، پوری \ealkeytr{صورت} کا \ealkeytr{عامل} نہیں، اس لیے اسے \ealkeytr{حذف} نہیں کیا جا سکتا۔}{\textbf{False.} $x^{2}$ is a \ealkey{term}, not a \ealkey{factor} of the whole \ealkey{numerator}, so it cannot be \ealkey{cancelled}.}
   $\dfrac{x(x+7)}{x(x+2)}=\dfrac{x+7}{x+2}$}

\tf[\LARGE][\large]{$\dfrac{x-3}{3-x}=1$}
  {\ealpara{\textbf{غلط۔} دونوں \ealkeytr{بریکٹ} ایک جیسے نہیں ہیں: $3-x=-(x-3)$۔}{\textbf{False.} The two \ealkey{brackets} are not the same: $3-x=-(x-3)$.}
   $\dfrac{x-3}{-(x-3)}=-1$
   \ealpara{جب $x=5$: $\dfrac{2}{-2}=-1$۔}{When $x=5$: $\dfrac{2}{-2}=-1$.}}

\tf[\LARGE][\large]{$\dfrac{x^{2}-25}{x-5}=x+5$}
  {\ealpara{\textbf{صحیح۔} \ealkeytr{صورت} \ealkeytr{دو مربعوں کا فرق} ہے: $x^{2}-25=(x-5)(x+5)$۔}{\textbf{True.} The \ealkey{numerator} is a \ealkey{difference of two squares}: $x^{2}-25=(x-5)(x+5)$.}
   $\dfrac{\strike{(x-5)}(x+5)}{\strike{(x-5)}}=x+5$}

\tf[\LARGE][\large]{$\dfrac{x^{2}-16}{x^{2}+8x+16}=1$}
  {\ealpara{\textbf{غلط۔} $(x-4)$ اور $(x+4)$ مختلف \ealkeytr{بریکٹ} ہیں اور \ealkeytr{حذف} نہیں ہوتے۔}{\textbf{False.} $(x-4)$ and $(x+4)$ are different \ealkey{brackets} and do not \ealkey{cancel}.}
   $\dfrac{(x-4)(x+4)}{(x+4)(x+4)}=\dfrac{x-4}{x+4}$}

\qq[\LARGE][\large]{\ealpara{\ealkeytr{سادہ} کریں $\dfrac{x^{2}+8x+15}{x+5}$}{\ealkey{Simplify} $\dfrac{x^{2}+8x+15}{x+5}$}}
  {\ealpara{\ealkeytr{تجزیه} کریں \ealkeytr{صورت} کو۔ دو اعداد جن کا \ealkeytr{حاصل جمع} $8$ اور \ealkeytr{حاصل ضرب} $15$ ہو: $3$ اور $5$۔}{\ealkey{Factorise} the \ealkey{numerator}. Two numbers with a \ealkey{sum} of $8$ and a \ealkey{product} of $15$: $3$ and $5$.}
   $\dfrac{(x+3)\strike{(x+5)}}{\strike{(x+5)}}=x+3$}

\qq[\LARGE][\large]{\ealpara{\ealkeytr{سادہ} کریں $\dfrac{x^{2}+7x+12}{x+3}$}{\ealkey{Simplify} $\dfrac{x^{2}+7x+12}{x+3}$}}
  {\ealpara{دو اعداد جن کا \ealkeytr{حاصل جمع} $7$ اور \ealkeytr{حاصل ضرب} $12$ ہو: $3$ اور $4$۔}{Two numbers with a \ealkey{sum} of $7$ and a \ealkey{product} of $12$: $3$ and $4$.}
   $\dfrac{\strike{(x+3)}(x+4)}{\strike{(x+3)}}=x+4$}

\qq[\LARGE][\large]{\ealpara{\ealkeytr{سادہ} کریں $\dfrac{x^{2}-36}{x+6}$}{\ealkey{Simplify} $\dfrac{x^{2}-36}{x+6}$}}
  {\ealpara{\ealkeytr{صورت} \ealkeytr{دو مربعوں کا فرق} ہے: $x^{2}-36=(x+6)(x-6)$۔}{The \ealkey{numerator} is a \ealkey{difference of two squares}: $x^{2}-36=(x+6)(x-6)$.}
   $\dfrac{\strike{(x+6)}(x-6)}{\strike{(x+6)}}=x-6$}

\qq[\LARGE][\normalsize]{\ealpara{\ealkeytr{سادہ} کریں $\dfrac{x^{2}+2x-8}{x-2}$}{\ealkey{Simplify} $\dfrac{x^{2}+2x-8}{x-2}$}}
  {\ealpara{دو اعداد جن کا \ealkeytr{حاصل جمع} $2$ اور \ealkeytr{حاصل ضرب} $-8$ ہو: $4$ اور $-2$۔}{Two numbers with a \ealkey{sum} of $2$ and a \ealkey{product} of $-8$: $4$ and $-2$.}
   $\dfrac{(x+4)\strike{(x-2)}}{\strike{(x-2)}}=x+4$
   \par\vspace{0.5em}
   \quadgrid{x}{4}{x}{-2}{x^{2}}{4x}{-2x}{-8}}

\qq[\LARGE][\large]{\ealpara{\ealkeytr{سادہ} کریں $\dfrac{x^{2}+6x+8}{x^{2}+5x+4}$}{\ealkey{Simplify} $\dfrac{x^{2}+6x+8}{x^{2}+5x+4}$}}
  {\ealpara{\ealkeytr{تجزیه} کریں \ealkeytr{صورت} اور \ealkeytr{مخرج} کو۔}{\ealkey{Factorise} the \ealkey{numerator} and the \ealkey{denominator}.}
   \par\vspace{0.4em}
   \work{$\dfrac{x^{2}+6x+8}{x^{2}+5x+4}$ & $\dfrac{(x+2)(x+4)}{(x+1)(x+4)}$\\[0.95em]
         $\dfrac{(x+2)\strike{(x+4)}}{(x+1)\strike{(x+4)}}$ & $\dfrac{x+2}{x+1}$}}

\qq[\LARGE][\large]{\ealpara{\ealkeytr{سادہ} کریں $\dfrac{x^{2}-x-6}{x^{2}-9}$}{\ealkey{Simplify} $\dfrac{x^{2}-x-6}{x^{2}-9}$}}
  {\ealpara{\ealkeytr{مخرج} \ealkeytr{دو مربعوں کا فرق} ہے۔}{The \ealkey{denominator} is a \ealkey{difference of two squares}.}
   \par\vspace{0.4em}
   \work{$\dfrac{x^{2}-x-6}{x^{2}-9}$ & $\dfrac{(x-3)(x+2)}{(x-3)(x+3)}$\\[0.95em]
         $\dfrac{\strike{(x-3)}(x+2)}{\strike{(x-3)}(x+3)}$ & $\dfrac{x+2}{x+3}$}}

\qq[\LARGE][\large]{\ealpara{\ealkeytr{سادہ} کریں $\dfrac{3x+6}{x^{2}+5x+6}$}{\ealkey{Simplify} $\dfrac{3x+6}{x^{2}+5x+6}$}}
  {\ealpara{\ealkeytr{صورت} کا \ealkeytr{مشترک عامل} $3$ ہے؛ \ealkeytr{مخرج} دو \ealkeytr{بریکٹوں} میں \ealkeytr{تجزیه} ہوتا ہے۔}{The \ealkey{numerator} has a \ealkey{common factor} of $3$; the \ealkey{denominator} \ealkey{factorises} into two \ealkey{brackets}.}
   \par\vspace{0.4em}
   \work{$\dfrac{3x+6}{x^{2}+5x+6}$ & $\dfrac{3(x+2)}{(x+2)(x+3)}$\\[0.95em]
         $\dfrac{3\strike{(x+2)}}{\strike{(x+2)}(x+3)}$ & $\dfrac{3}{x+3}$}}

\qq[\LARGE][\normalsize]{\ealpara{\ealkeytr{سادہ} کریں $\dfrac{2x^{2}+7x+3}{x+3}$}{\ealkey{Simplify} $\dfrac{2x^{2}+7x+3}{x+3}$}}
  {\ealpara{$x^{2}$ کا \ealkeytr{عددی سر} $2$ ہے، اس لیے ایک \ealkeytr{بریکٹ} $2x$ سے شروع ہوتا ہے۔}{The \ealkey{coefficient} of $x^{2}$ is $2$, so one \ealkey{bracket} starts with $2x$.}
   $\dfrac{(2x+1)\strike{(x+3)}}{\strike{(x+3)}}=2x+1$
   \par\vspace{0.5em}
   \quadgrid{2x}{1}{x}{3}{2x^{2}}{x}{6x}{3}}

\qq[\LARGE][\large]{\ealpara{\ealkeytr{سادہ} کریں $\dfrac{2x^{2}-8}{x^{2}+4x+4}$}{\ealkey{Simplify} $\dfrac{2x^{2}-8}{x^{2}+4x+4}$}}
  {\ealpara{پہلے \ealkeytr{صورت} میں سے $2$ کا \ealkeytr{مشترک عامل} باہر نکالیں۔}{Take the \ealkey{common factor} of $2$ out of the \ealkey{numerator} first.}
   \par\vspace{0.4em}
   \work{$\dfrac{2x^{2}-8}{x^{2}+4x+4}$ & $\dfrac{2(x^{2}-4)}{(x+2)(x+2)}$\\[0.95em]
         $\dfrac{2(x^{2}-4)}{(x+2)(x+2)}$ & $\dfrac{2(x-2)\strike{(x+2)}}{(x+2)\strike{(x+2)}}$\\[0.95em]
         $\dfrac{2x^{2}-8}{x^{2}+4x+4}$ & $\dfrac{2(x-2)}{x+2}$}}

\qq[\LARGE][\large]{\ealpara{\ealkeytr{سادہ} کریں $\dfrac{6x^{2}+x-2}{4x^{2}-1}$}{\ealkey{Simplify} $\dfrac{6x^{2}+x-2}{4x^{2}-1}$}}
  {\ealpara{\ealkeytr{مخرج} \ealkeytr{دو مربعوں کا فرق} ہے۔}{The \ealkey{denominator} is a \ealkey{difference of two squares}.}
   \par\vspace{0.4em}
   \work{$6x^{2}+x-2$ & $(3x+2)(2x-1)$\\[0.95em]
         $4x^{2}-1$ & $(2x-1)(2x+1)$\\[0.95em]
         $\dfrac{(3x+2)\strike{(2x-1)}}{(2x+1)\strike{(2x-1)}}$ & $\dfrac{3x+2}{2x+1}$}}

%================================================================
\qtopic{\ealkey{Adding} and \ealkey{subtracting} \ealkey{algebraic} \ealkey{fractions}}
%================================================================

\tf[\LARGE][\Large]{$\dfrac{2}{x}+\dfrac{3}{x}=\dfrac{5}{2x}$}
  {\ealpara{\textbf{غلط۔} \ealkeytr{مخرج} پہلے سے ایک جیسے ہیں، اس لیے صرف \ealkeytr{صورتیں} \ealkeytr{جمع} کی جاتی ہیں۔}{\textbf{False.} The \ealkey{denominators} are already the same, so only the \ealkey{numerators} are \ealkey{added}.}
   $\dfrac{2}{x}+\dfrac{3}{x}=\dfrac{5}{x}$}

\tf[\LARGE][\Large]{$\dfrac{x}{3}+\dfrac{x}{4}=\dfrac{7x}{12}$}
  {\ealpara{\textbf{صحیح۔} \ealkeytr{مشترک مخرج} $12$ ہے۔}{\textbf{True.} The \ealkey{common denominator} is $12$.}
   $\dfrac{4x}{12}+\dfrac{3x}{12}=\dfrac{7x}{12}$}

\tf[\LARGE][\large]{$\dfrac{1}{x}+\dfrac{1}{y}=\dfrac{1}{x+y}$}
  {\ealpara{\textbf{غلط۔} اسے $x=2$ اور $y=3$ کے ساتھ آزمائیں: بائیں طرف $\dfrac{5}{6}$ ہے اور دائیں طرف $\dfrac{1}{5}$ ہے۔}{\textbf{False.} Test it with $x=2$ and $y=3$: the left side is $\dfrac{5}{6}$ and the right side is $\dfrac{1}{5}$.}
   \ealpara{\ealkeytr{مشترک مخرج} $xy$ ہے، اس لیے جواب $\dfrac{x+y}{xy}$ ہے۔}{The \ealkey{common denominator} is $xy$, so the answer is $\dfrac{x+y}{xy}$.}}

\tf[\Large][\large]{$\dfrac{4}{x}-\dfrac{x-3}{x}=\dfrac{4-x-3}{x}$}
  {\ealpara{\textbf{غلط۔} منفی \ealkeytr{علامت} پوری دوسری \ealkeytr{صورت} پر لاگو ہوتی ہے، اس لیے $-3$، $+3$ بن جاتا ہے۔}{\textbf{False.} The minus \ealkey{sign} applies to the whole second \ealkey{numerator}, so the $-3$ becomes $+3$.}
   $\dfrac{4-(x-3)}{x}=\dfrac{4-x+3}{x}=\dfrac{7-x}{x}$}

\tf[\Large][\large]{\ealpara{$\dfrac{1}{x+1}$ اور $\dfrac{2}{x+3}$ کا \ealkeytr{مشترک مخرج} $(x+1)(x+3)$ ہے}{The \ealkey{common denominator} of $\dfrac{1}{x+1}$ and $\dfrac{2}{x+3}$ is $(x+1)(x+3)$}}
  {\ealpara{\textbf{صحیح۔} \ealkeytr{مخرج} کو \ealkeytr{ضرب} دیا جاتا ہے، \ealkeytr{جمع} نہیں کیا جاتا۔}{\textbf{True.} The \ealkey{denominators} are \ealkey{multiplied}, not \ealkey{added}.}
   \ealpara{\ealkeytr{مشترک مخرج} $(x+1)(x+3)$ ہے، $x+4$ نہیں۔}{The \ealkey{common denominator} is $(x+1)(x+3)$, not $x+4$.}}

\qq[\LARGE][\large]{\ealpara{ایک واحد \ealkeytr{کسر} کے طور پر لکھیں\\[0.4em] $\dfrac{3}{x}+\dfrac{4}{x}$}{Write as a single \ealkey{fraction}\\[0.4em] $\dfrac{3}{x}+\dfrac{4}{x}$}}
  {\ealpara{\ealkeytr{مخرج} ایک جیسے ہیں، اس لیے \ealkeytr{صورتیں} \ealkeytr{جمع} کی جاتی ہیں۔}{The \ealkey{denominators} are the same, so the \ealkey{numerators} are \ealkey{added}.}
   $\dfrac{3}{x}+\dfrac{4}{x}=\dfrac{7}{x}$}

\qq[\LARGE][\large]{\ealpara{ایک واحد \ealkeytr{کسر} کے طور پر لکھیں\\[0.4em] $\dfrac{x}{2}+\dfrac{x}{5}$}{Write as a single \ealkey{fraction}\\[0.4em] $\dfrac{x}{2}+\dfrac{x}{5}$}}
  {\ealpara{\ealkeytr{مشترک مخرج} $10$ ہے۔}{The \ealkey{common denominator} is $10$.}
   \par\vspace{0.4em}
   \work{$\dfrac{x}{2}+\dfrac{x}{5}$ & $\dfrac{5x}{10}+\dfrac{2x}{10}$\\[0.95em]
         $\dfrac{5x}{10}+\dfrac{2x}{10}$ & $\dfrac{7x}{10}$}}

\qq[\LARGE][\large]{\ealpara{ایک واحد \ealkeytr{کسر} کے طور پر لکھیں\\[0.4em] $\dfrac{x}{4}+\dfrac{3x}{8}$}{Write as a single \ealkey{fraction}\\[0.4em] $\dfrac{x}{4}+\dfrac{3x}{8}$}}
  {\ealpara{$8$، $4$ کا \ealkeytr{مضاعف} ہے، اس لیے \ealkeytr{مشترک مخرج} $8$ ہے۔}{$8$ is a \ealkey{multiple} of $4$, so the \ealkey{common denominator} is $8$.}
   \par\vspace{0.4em}
   \work{$\dfrac{x}{4}+\dfrac{3x}{8}$ & $\dfrac{2x}{8}+\dfrac{3x}{8}$\\[0.95em]
         $\dfrac{2x}{8}+\dfrac{3x}{8}$ & $\dfrac{5x}{8}$}}

\qq[\LARGE][\large]{\ealpara{ایک واحد \ealkeytr{کسر} کے طور پر لکھیں\\[0.4em] $\dfrac{2x}{5}-\dfrac{x}{3}$}{Write as a single \ealkey{fraction}\\[0.4em] $\dfrac{2x}{5}-\dfrac{x}{3}$}}
  {\ealpara{\ealkeytr{مشترک مخرج} $15$ ہے۔}{The \ealkey{common denominator} is $15$.}
   \par\vspace{0.4em}
   \work{$\dfrac{2x}{5}-\dfrac{x}{3}$ & $\dfrac{6x}{15}-\dfrac{5x}{15}$\\[0.95em]
         $\dfrac{6x}{15}-\dfrac{5x}{15}$ & $\dfrac{x}{15}$}}

\qq[\LARGE][\large]{\ealpara{ایک واحد \ealkeytr{کسر} کے طور پر لکھیں\\[0.4em] $\dfrac{x+1}{2}+\dfrac{x+2}{3}$}{Write as a single \ealkey{fraction}\\[0.4em] $\dfrac{x+1}{2}+\dfrac{x+2}{3}$}}
  {\ealpara{\ealkeytr{مشترک مخرج} $6$ ہے۔ ہر \ealkeytr{صورت} کو اسی \ealkeytr{عامل} سے \ealkeytr{ضرب} دیا جاتا ہے جس سے اس کا \ealkeytr{مخرج} ضرب دیا گیا ہو۔}{The \ealkey{common denominator} is $6$. Each \ealkey{numerator} is \ealkey{multiplied} by the same \ealkey{factor} as its \ealkey{denominator}.}
   \par\vspace{0.4em}
   \work{$\dfrac{x+1}{2}+\dfrac{x+2}{3}$ & $\dfrac{3(x+1)+2(x+2)}{6}$\\[0.95em]
         $\dfrac{3x+3+2x+4}{6}$ & $\dfrac{5x+7}{6}$}}

\qq[\LARGE][\large]{\ealpara{ایک واحد \ealkeytr{کسر} کے طور پر لکھیں\\[0.4em] $\dfrac{2x+3}{4}-\dfrac{x-1}{6}$}{Write as a single \ealkey{fraction}\\[0.4em] $\dfrac{2x+3}{4}-\dfrac{x-1}{6}$}}
  {\ealpara{\ealkeytr{مشترک مخرج} $12$ ہے۔ منفی \ealkeytr{علامت} پوری دوسری \ealkeytr{صورت} پر لاگو ہوتی ہے۔}{The \ealkey{common denominator} is $12$. The minus \ealkey{sign} applies to the whole second \ealkey{numerator}.}
   \par\vspace{0.4em}
   \work{$\dfrac{2x+3}{4}-\dfrac{x-1}{6}$ & $\dfrac{3(2x+3)-2(x-1)}{12}$\\[0.95em]
         $\dfrac{6x+9-2x+2}{12}$ & $\dfrac{4x+11}{12}$}}

\qq[\LARGE][\large]{\ealpara{ایک واحد \ealkeytr{کسر} کے طور پر لکھیں\\[0.4em] $\dfrac{2}{x}+\dfrac{3}{y}$}{Write as a single \ealkey{fraction}\\[0.4em] $\dfrac{2}{x}+\dfrac{3}{y}$}}
  {\ealpara{\ealkeytr{مشترک مخرج} $xy$ ہے۔}{The \ealkey{common denominator} is $xy$.}
   \par\vspace{0.4em}
   \work{$\dfrac{2}{x}+\dfrac{3}{y}$ & $\dfrac{2y}{xy}+\dfrac{3x}{xy}$\\[0.95em]
         $\dfrac{2y}{xy}+\dfrac{3x}{xy}$ & $\dfrac{3x+2y}{xy}$}}

\qq[\LARGE][\large]{\ealpara{ایک واحد \ealkeytr{کسر} کے طور پر لکھیں\\[0.4em] $\dfrac{1}{x+1}+\dfrac{2}{x+3}$}{Write as a single \ealkey{fraction}\\[0.4em] $\dfrac{1}{x+1}+\dfrac{2}{x+3}$}}
  {\ealpara{\ealkeytr{مشترک مخرج} $(x+1)(x+3)$ ہے۔}{The \ealkey{common denominator} is $(x+1)(x+3)$.}
   \par\vspace{0.4em}
   \work{$\dfrac{1}{x+1}+\dfrac{2}{x+3}$ & $\dfrac{(x+3)+2(x+1)}{(x+1)(x+3)}$\\[0.95em]
         $\dfrac{x+3+2x+2}{(x+1)(x+3)}$ & $\dfrac{3x+5}{(x+1)(x+3)}$}}

\qq[\LARGE][\large]{\ealpara{ایک واحد \ealkeytr{کسر} کے طور پر لکھیں\\[0.4em] $\dfrac{3}{x-1}-\dfrac{2}{x+2}$}{Write as a single \ealkey{fraction}\\[0.4em] $\dfrac{3}{x-1}-\dfrac{2}{x+2}$}}
  {\ealpara{\ealkeytr{مشترک مخرج} $(x-1)(x+2)$ ہے۔ منفی \ealkeytr{علامت} پوری دوسری \ealkeytr{صورت} پر لاگو ہوتی ہے۔}{The \ealkey{common denominator} is $(x-1)(x+2)$. The minus \ealkey{sign} applies to the whole second \ealkey{numerator}.}
   \par\vspace{0.4em}
   \work{$\dfrac{3}{x-1}-\dfrac{2}{x+2}$ & $\dfrac{3(x+2)-2(x-1)}{(x-1)(x+2)}$\\[0.95em]
         $\dfrac{3x+6-2x+2}{(x-1)(x+2)}$ & $\dfrac{x+8}{(x-1)(x+2)}$}}

\qq[\LARGE][\small]{\ealpara{ایک واحد \ealkeytr{کسر} کے طور پر لکھیں\\[0.4em] $\dfrac{1}{x-2}+\dfrac{3}{x^{2}-4}$}{Write as a single \ealkey{fraction}\\[0.4em] $\dfrac{1}{x-2}+\dfrac{3}{x^{2}-4}$}}
  {\ealpara{\ealkeytr{تجزیه} کریں دوسرے \ealkeytr{مخرج} کو، پھر دونوں \ealkeytr{کسروں} کو $(x-2)(x+2)$ کے اوپر لکھیں۔}{\ealkey{Factorise} the second \ealkey{denominator}, then write both \ealkey{fractions} over $(x-2)(x+2)$.}
   \par\vspace{0.4em}
   \work{$x^{2}-4$ & $(x-2)(x+2)$\\[1.3em]
         $\dfrac{1}{x-2}$ & $\dfrac{x+2}{(x-2)(x+2)}$\\[1.3em]
         $\dfrac{1}{x-2}+\dfrac{3}{x^{2}-4}$ & $\dfrac{x+2+3}{(x-2)(x+2)}$\\[1.3em]
         $\dfrac{x+2+3}{(x-2)(x+2)}$ & $\dfrac{x+5}{(x-2)(x+2)}$}}

%================================================================
\qtopic{\ealkey{Multiplying} and \ealkey{dividing} \ealkey{algebraic} \ealkey{fractions}}
%================================================================

\tf[\Large][\large]{\ealpara{$\dfrac{a}{b}\div\dfrac{c}{d}$ معلوم کرنے کے لیے آپ پہلی \ealkeytr{کسر} کو الٹ دیتے ہیں}{To work out $\dfrac{a}{b}\div\dfrac{c}{d}$ you turn the first \ealkey{fraction} upside down}}
  {\ealpara{\textbf{غلط۔} دوسری \ealkeytr{کسر} کو الٹا جاتا ہے، اور پھر دونوں \ealkeytr{کسروں} کو \ealkeytr{ضرب} دیا جاتا ہے۔}{\textbf{False.} It is the second \ealkey{fraction} that is turned upside down, and the two \ealkey{fractions} are then \ealkey{multiplied}.}
   $\dfrac{a}{b}\div\dfrac{c}{d}=\dfrac{a}{b}\times\dfrac{d}{c}$}

\tf[\Large][\large]{\ealpara{$\dfrac{x}{2}\times\dfrac{3}{5}$ کو \ealkeytr{مشترک مخرج} معلوم کیے بغیر حل کیا جا سکتا ہے}{$\dfrac{x}{2}\times\dfrac{3}{5}$ can be worked out without finding a \ealkey{common denominator}}}
  {\ealpara{\textbf{صحیح۔} \ealkeytr{مشترک مخرج} \ealkeytr{جمع} اور \ealkeytr{تفریق} کے لیے درکار ہوتا ہے، \ealkeytr{ضرب} کے لیے نہیں۔}{\textbf{True.} A \ealkey{common denominator} is needed for \ealkey{adding} and \ealkey{subtracting}, not for \ealkey{multiplying}.}
   \ealpara{\ealkeytr{صورتیں} \ealkeytr{ضرب} دی جاتی ہیں اور \ealkeytr{مخرج} \ealkeytr{ضرب} دیے جاتے ہیں، اس لیے جواب $\dfrac{3x}{10}$ ہے۔}{The \ealkey{numerators} are \ealkey{multiplied} and the \ealkey{denominators} are \ealkey{multiplied}, so the answer is $\dfrac{3x}{10}$.}}

\tf[\LARGE][\large]{$\dfrac{a}{b}\times\dfrac{c}{d}=\dfrac{a+c}{b+d}$}
  {\ealpara{\textbf{غلط۔} اسے $\dfrac{1}{2}\times\dfrac{1}{2}$ کے ساتھ آزمائیں: بائیں طرف $\dfrac{1}{4}$ ہے اور دائیں طرف $\dfrac{2}{4}=\dfrac{1}{2}$ ہے۔}{\textbf{False.} Test it with $\dfrac{1}{2}\times\dfrac{1}{2}$: the left side is $\dfrac{1}{4}$ and the right side is $\dfrac{2}{4}=\dfrac{1}{2}$.}
   $\dfrac{a}{b}\times\dfrac{c}{d}=\dfrac{ac}{bd}$}

\tf[\Large][\large]{$\dfrac{x}{4}\div 2=\dfrac{x}{2}$}
  {\ealpara{\textbf{غلط۔} پورا عدد $2$ \ealkeytr{کسر} $\dfrac{2}{1}$ ہے، اس لیے اسے الٹا جاتا ہے اور دونوں کو \ealkeytr{ضرب} دیا جاتا ہے۔}{\textbf{False.} The whole number $2$ is the \ealkey{fraction} $\dfrac{2}{1}$, so it is turned upside down and the two are \ealkey{multiplied}.}
   $\dfrac{x}{4}\div\dfrac{2}{1}=\dfrac{x}{4}\times\dfrac{1}{2}=\dfrac{x}{8}$
   \ealpara{جب $x=8$: $2\div 2=1$، نہ کہ $4$۔}{When $x=8$: $2\div 2=1$, not $4$.}}

\tf[\Large][\large]{\ealpara{$\dfrac{x+1}{4}\times\dfrac{8}{x+1}$ \ealkeytr{سادہ} ہو کر $2$ بنتا ہے}{$\dfrac{x+1}{4}\times\dfrac{8}{x+1}$ \ealkey{simplifies} to $2$}}
  {\ealpara{\textbf{صحیح۔} ضرب کی \ealkeytr{علامت} کے آر پار \ealkeytr{حذف} کرنا جائز ہے، کیونکہ $x+1$ پوری \ealkeytr{صورت} اور پورے \ealkeytr{مخرج} کا \ealkeytr{عامل} ہے۔}{\textbf{True.} \ealkey{Cancelling} across a multiplication \ealkey{sign} is allowed, because $x+1$ is a \ealkey{factor} of a whole \ealkey{numerator} and of a whole \ealkey{denominator}.}
   $\dfrac{8\strike{(x+1)}}{4\strike{(x+1)}}=\dfrac{8}{4}=2$}

\qq[\LARGE][\large]{\ealpara{\ealkeytr{سادہ} کریں\\[0.4em] $\dfrac{x}{2}\times\dfrac{3}{x}$}{\ealkey{Simplify}\\[0.4em] $\dfrac{x}{2}\times\dfrac{3}{x}$}}
  {\ealpara{\ealkeytr{ضرب} دیں \ealkeytr{صورتوں} کو اور \ealkeytr{ضرب} دیں \ealkeytr{مخرجوں} کو، پھر \ealkeytr{مشترک عامل} $x$ کو \ealkeytr{حذف} کریں۔}{\ealkey{Multiply} the \ealkey{numerators} and \ealkey{multiply} the \ealkey{denominators}, then \ealkey{cancel} the \ealkey{common factor} $x$.}
   \par\vspace{0.4em}
   \work{$\dfrac{x}{2}\times\dfrac{3}{x}$ & $\dfrac{3x}{2x}$\\[0.95em]
         $\dfrac{3\strike{x}}{2\strike{x}}$ & $\dfrac{3}{2}$}}

\qq[\LARGE][\large]{\ealpara{\ealkeytr{سادہ} کریں\\[0.4em] $\dfrac{3x}{4}\times\dfrac{8}{9x}$}{\ealkey{Simplify}\\[0.4em] $\dfrac{3x}{4}\times\dfrac{8}{9x}$}}
  {\work{$\dfrac{3x}{4}\times\dfrac{8}{9x}$ & $\dfrac{24x}{36x}$\\[0.95em]
         $\dfrac{24\strike{x}}{36\strike{x}}$ & $\dfrac{24}{36}$\\[0.95em]
         $\dfrac{24}{36}$ & $\dfrac{2}{3}$}}

\qq[\LARGE][\large]{\ealpara{\ealkeytr{سادہ} کریں\\[0.4em] $\dfrac{x}{5}\div\dfrac{x}{2}$}{\ealkey{Simplify}\\[0.4em] $\dfrac{x}{5}\div\dfrac{x}{2}$}}
  {\ealpara{دوسری \ealkeytr{کسر} کو الٹ دیں اور \ealkeytr{ضرب} دیں}{turn the second \ealkey{fraction} upside down and \ealkey{multiply}}
   \ealpara{\ealkeytr{مشترک عامل} $x$ کو \ealkeytr{حذف} کریں}{\ealkey{cancel} the \ealkey{common factor} $x$}
   \flow{\tfrac{x}{5}\div\tfrac{x}{2}}{turn the second fraction upside
   down and multiply}{\tfrac{x}{5}\times\tfrac{2}{x}}{cancel the common
   factor $x$}{\tfrac{2}{5}}}

\qq[\LARGE][\large]{\ealpara{\ealkeytr{سادہ} کریں\\[0.4em] $\dfrac{x}{3}\div\dfrac{2}{x}$}{\ealkey{Simplify}\\[0.4em] $\dfrac{x}{3}\div\dfrac{2}{x}$}}
  {\ealpara{دوسری \ealkeytr{کسر} کو الٹ دیں اور \ealkeytr{ضرب} دیں}{turn the second \ealkey{fraction} upside down and \ealkey{multiply}}
   \ealpara{\ealkeytr{صورتوں} اور \ealkeytr{مخرجوں} کو \ealkeytr{ضرب} دیں}{\ealkey{multiply} the \ealkey{numerators} and the \ealkey{denominators}}
   \flow{\tfrac{x}{3}\div\tfrac{2}{x}}{turn the second fraction upside
   down and multiply}{\tfrac{x}{3}\times\tfrac{x}{2}}{multiply the
   numerators and the denominators}{\tfrac{x^{2}}{6}}}

\qq[\LARGE][\large]{\ealpara{\ealkeytr{سادہ} کریں\\[0.4em] $\dfrac{x+1}{3}\times\dfrac{6}{x+2}$}{\ealkey{Simplify}\\[0.4em] $\dfrac{x+1}{3}\times\dfrac{6}{x+2}$}}
  {\ealpara{دونوں \ealkeytr{بریکٹ} مختلف ہیں، اس لیے صرف اعداد \ealkeytr{حذف} ہوتے ہیں۔}{The two \ealkey{brackets} are different, so only the numbers \ealkey{cancel}.}
   \par\vspace{0.4em}
   \work{$\dfrac{x+1}{3}\times\dfrac{6}{x+2}$ & $\dfrac{6(x+1)}{3(x+2)}$\\[0.95em]
         $\dfrac{6(x+1)}{3(x+2)}$ & $\dfrac{2(x+1)}{x+2}$}}

\qq[\LARGE][\small]{\ealpara{\ealkeytr{سادہ} کریں\\[0.4em] $\dfrac{x+2}{5}\div\dfrac{x+2}{10}$}{\ealkey{Simplify}\\[0.4em] $\dfrac{x+2}{5}\div\dfrac{x+2}{10}$}}
  {\ealpara{دوسری \ealkeytr{کسر} کو الٹ دیں اور \ealkeytr{ضرب} دیں۔}{Turn the second \ealkey{fraction} upside down and \ealkey{multiply}.}
   \par\vspace{0.4em}
   \work{$\dfrac{x+2}{5}\div\dfrac{x+2}{10}$ & $\dfrac{x+2}{5}\times\dfrac{10}{x+2}$\\[0.95em]
         $\dfrac{x+2}{5}\times\dfrac{10}{x+2}$ & $\dfrac{10(x+2)}{5(x+2)}$\\[0.95em]
         $\dfrac{10\strike{(x+2)}}{5\strike{(x+2)}}$ & $2$}}

\qq[\LARGE][\large]{\ealpara{\ealkeytr{سادہ} کریں\\[0.4em] $\dfrac{5x^{3}}{4}\div\dfrac{x}{8}$}{\ealkey{Simplify}\\[0.4em] $\dfrac{5x^{3}}{4}\div\dfrac{x}{8}$}}
  {\ealpara{دوسری \ealkeytr{کسر} کو الٹ دیں اور \ealkeytr{ضرب} دیں}{turn the second \ealkey{fraction} upside down and \ealkey{multiply}}
   \ealpara{اعداد کو \ealkeytr{تقسیم} کریں اور \ealkeytr{اس} کو \ealkeytr{تفریق} کریں}{\ealkey{divide} the numbers and \ealkey{subtract} the \ealkey{indices}}
   \flow{\tfrac{5x^{3}}{4}\div\tfrac{x}{8}}{turn the second fraction
   upside down and multiply}{\tfrac{5x^{3}}{4}\times\tfrac{8}{x}}{divide
   the numbers and subtract the indices}{10x^{2}}}

\qq[\Large][\small]{\ealpara{\ealkeytr{سادہ} کریں\\[0.4em] $\dfrac{x^{2}-9}{x+1}\times\dfrac{x+1}{x-3}$}{\ealkey{Simplify}\\[0.4em] $\dfrac{x^{2}-9}{x+1}\times\dfrac{x+1}{x-3}$}}
  {\ealpara{$x^{2}-9$ \ealkeytr{دو مربعوں کا فرق} ہے۔}{$x^{2}-9$ is the \ealkey{difference of two squares}.}
   \par\vspace{0.4em}
   \work{$x^{2}-9$ & $(x-3)(x+3)$\\[0.95em]
         $\dfrac{x^{2}-9}{x+1}\times\dfrac{x+1}{x-3}$ & $\dfrac{(x-3)(x+3)(x+1)}{(x+1)(x-3)}$\\[0.95em]
         $\dfrac{\strike{(x-3)}(x+3)\strike{(x+1)}}{\strike{(x+1)}\strike{(x-3)}}$ & $x+3$}}

\qq[\Large][\small]{\ealpara{\ealkeytr{سادہ} کریں\\[0.4em] $\dfrac{x^{2}+7x+10}{x-1}\div\dfrac{x+2}{x-1}$}{\ealkey{Simplify}\\[0.4em] $\dfrac{x^{2}+7x+10}{x-1}\div\dfrac{x+2}{x-1}$}}
  {\ealpara{\ealkeytr{تجزیه} کریں \ealkeytr{دو درجی} کو، پھر دوسری \ealkeytr{کسر} کو الٹ دیں اور \ealkeytr{ضرب} دیں۔}{\ealkey{Factorise} the \ealkey{quadratic}, then turn the second \ealkey{fraction} upside down and \ealkey{multiply}.}
   \par\vspace{0.4em}
   \work{$x^{2}+7x+10$ & $(x+2)(x+5)$\\[0.95em]
         $\dfrac{(x+2)(x+5)}{x-1}\div\dfrac{x+2}{x-1}$ & $\dfrac{(x+2)(x+5)}{x-1}\times\dfrac{x-1}{x+2}$\\[0.95em]
         $\dfrac{\strike{(x+2)}(x+5)\strike{(x-1)}}{\strike{(x-1)}\strike{(x+2)}}$ & $x+5$}}

\qq[\Large][\footnotesize]{\ealpara{\ealkeytr{سادہ} کریں\\[0.4em] $\dfrac{x^{2}+3x+2}{x^{2}-1}\div\dfrac{x^{2}+4x+4}{x^{2}-x}$}{\ealkey{Simplify}\\[0.4em] $\dfrac{x^{2}+3x+2}{x^{2}-1}\div\dfrac{x^{2}+4x+4}{x^{2}-x}$}}
  {\ealpara{\ealkeytr{تجزیه} کریں تمام چار \ealkeytr{دو درجی} اظہارات کو، پھر دوسری \ealkeytr{کسر} کو الٹ دیں اور \ealkeytr{ضرب} دیں۔}{\ealkey{Factorise} all four \ealkey{quadratics}, then turn the second \ealkey{fraction} upside down and \ealkey{multiply}.}
   \par\vspace{0.4em}
   \work{$x^{2}+3x+2$ & $(x+1)(x+2)$\\[0.95em]
         $x^{2}-1$ & $(x+1)(x-1)$\\[0.95em]
         $x^{2}+4x+4$ & $(x+2)(x+2)$\\[0.95em]
         $x^{2}-x$ & $x(x-1)$\\[0.95em]
         $\dfrac{(x+1)(x+2)}{(x+1)(x-1)}\times\dfrac{x(x-1)}{(x+2)(x+2)}$ & $\dfrac{(x+1)(x+2)\,x\,(x-1)}{(x+1)(x-1)(x+2)(x+2)}$\\[1.5em]
         $\dfrac{\strike{(x+1)}\strike{(x+2)}\,x\,\strike{(x-1)}}{\strike{(x+1)}\strike{(x-1)}\strike{(x+2)}(x+2)}$ & $\dfrac{x}{x+2}$}}

\end{document}